% BOOK-FILE-ID: f-94a6f53f3fe9
\section{用語一覧}\label{節：用語一覧}
本文で使う用語と,その説明がある場所をまとめる.
\paragraphbf{数列と漸化式}
\begin{description}[style=nextline,leftmargin=0pt,labelsep=0pt,itemsep=5pt]
  \item[\hypertarget{term:recurrence}{漸化式}] 項どうしの関係によって数列を定める式. \booknavref{節：数列と漸化式}{説明へ.}
  \item[\hypertarget{term:sequence}{数列}] 数を順に並べたもの.添字で各項の位置を表す. \booknavref{節：数列と漸化式}{説明へ.}
  \item[\hypertarget{term:initial}{初期条件}] 数列を作り始めるために与える値. \booknavref{節：数列と漸化式}{説明へ.}
  \item[\hypertarget{term:general}{一般項}] 番号からその位置の項を直接表す式. \booknavref{節：数列と漸化式}{説明へ.}
  \item[\hypertarget{term:arithmetic}{等差数列}] 次の項へ進むとき,同じ数を足す数列. \booknavref{節：等差型}{説明へ.}
  \item[\hypertarget{term:geometric}{等比数列}] 次の項へ進むとき,同じ数を掛ける数列. \booknavref{節：等比型}{説明へ.}
  \item[\hypertarget{term:difference}{階差}] 隣り合う項の差$a_{n+1}-a_n$のこと.階差の和に初項を加えると元の数列になる. \booknavref{節：階差型}{説明へ.}
  \item[\hypertarget{term:ratio}{階比}] 隣り合う項の比$\frac{a_{n+1}}{a_n}$のこと.分母が0でない範囲で考える. \booknavref{節：階比型}{説明へ.}
  \item[\hypertarget{term:sigma}{シグマ}] 和を,添字の範囲とともに書く記号. \booknavref{節：シグマの計算}{説明へ.}
  \item[\hypertarget{term:product}{総積}] 指定した範囲の項をすべて掛けたもの. \booknavref{節：階比型}{説明へ.}
  \item[\hypertarget{term:constant}{定数係数}] 未知の項に掛かる係数が,添字によらず一定であること. \booknavref{節：定数係数・変数係数}{説明へ.}
  \item[\hypertarget{term:variable}{変数係数}] 未知の項に掛かる係数が,添字によって変わること. \booknavref{節：定数係数・変数係数}{説明へ.}
  \item[\hypertarget{term:double}{二重数列}] 2つの添字で項を指定する数列. \booknavref{節：二方向の階差}{説明へ.}
  \item[\hypertarget{term:periodic}{周期性}] 一定の間隔を進めると同じ状態に戻る性質. \booknavref{節：非線形の周期性}{説明へ.}
  \item[\hypertarget{term:floor}{ガウス記号}] 実数$x$を超えない最大の整数を表す記号.\,$[x]$や$\lfloor x\rfloor$と書く. \booknavref{節：ガウス記号を含む式}{説明へ.}
  \item[\hypertarget{term:floor-function}{床関数}] ガウス記号$\lfloor x\rfloor$の大学以降での呼び方. \booknavref{節：ガウス記号を含む式}{説明へ.}
\end{description}
\paragraphbf{式を整える}
\begin{description}[style=nextline,leftmargin=0pt,labelsep=0pt,itemsep=5pt]
  \item[\hypertarget{term:auxiliary}{補助数列}] 元の数列を調べるために,\,$b_n=a_n-\alpha$のように新しく置いた数列. \booknavref{節：定数数列}{説明へ.}
  \item[\hypertarget{term:addend}{付加項}] 漸化式の右辺で,\,$a_n$などの未知の項を含まずに足されている項.\,$a_{n+1}=pa_n+f(n)$の$f(n)$など. \booknavref{節：階差等比複合型}{説明へ.}
  \item[\hypertarget{term:constant-sequence}{定数数列}] すべての項が同じ値をとる数列. \booknavref{節：定数数列}{説明へ.}
  \item[\hypertarget{term:normalize}{正規化}] 本書では,階比型の解$h_n$で両辺を割り,\,$a_n$の係数を1にする操作をこう呼ぶ. \booknavref{節：添字ごとに置き換え方を変える}{説明へ.}
  \item[\hypertarget{term:partial}{部分分数分解}] 1つの分数を,分母がより簡単な分数の和へ直す操作. \booknavref{節：部分分数分解}{説明へ.}
  \item[\hypertarget{term:square}{平方完成}] 二次式を,\,$a(x-p)^2+q$のように一次式の二乗を用いて表す操作. \booknavref{節：同じ数を足して引く}{説明へ.}
  \item[\hypertarget{term:fixed}{不動点}] $a_{n+1}=f(a_n)$で,\,$f(\alpha)=\alpha$を満たす値$\alpha$.その値から始めると数列は変化しない. \booknavref{節：特性方程式型}{説明へ.}
  \item[\hypertarget{term:fixed-function}{不動関数}] 本書独自の用語.階差等比複合型などで,元の漸化式を満たすように選んで引く関数.一般には特殊解と呼ばれる. \booknavref{節：階差等比複合型}{説明へ.}
  \item[\hypertarget{term:undetermined}{未定係数法}] 式の形を先に選び,係数を比較して未知の係数を決める方法. \booknavref{節：階差等比複合型}{説明へ.}
  \item[\hypertarget{term:induction}{数学的帰納法}] 出発点と,次の番号へ進めることを示して,すべての番号について証明する方法. \booknavref{節：数学的帰納法}{説明へ.}
  \item[\hypertarget{term:logarithm}{対数}] ある数が底の何乗かを表す数.真数と底の条件を伴う. \booknavref{節：対数}{説明へ.}
  \item[\hypertarget{term:inverse}{逆関数}] 関数の値から元の入力をただ1つに定める関数.元の関数の定義域と値域を入れ替えたものになる. \booknavref{節：対数}{説明へ.}
  \item[\hypertarget{term:binomial}{二項定理}] 2つの数の和の0以上の整数乗を,項の和へ展開する公式. \booknavref{節：二項定理}{説明へ.}
\end{description}
\paragraphbf{解の構造}
\begin{description}[style=nextline,leftmargin=0pt,labelsep=0pt,itemsep=5pt]
  \item[\hypertarget{term:linear}{線形}] 和と定数倍を保つ性質.線形漸化式とは,未知の項に係数を掛けて足した形だけで表される漸化式である. \booknavref{節：線形・非線形}{説明へ.}
  \item[\hypertarget{term:nonlinear}{非線形}] 線形の条件を満たさないもの.項の二乗や逆数を含む例を扱う. \booknavref{節：線形・非線形}{説明へ.}
  \item[\hypertarget{term:homogeneous}{同次}] 付加項のない線形漸化式のこと.斉次ともいう.付加項がある場合は非同次という. \booknavref{節：同次・重ね合わせ}{説明へ.}
  \item[\hypertarget{term:characteristic}{特性方程式}] 漸化式の解法に必要な値を求めるための方程式.特性方程式型・一次分数型では不動点を,隣接3項型・連立型では等比型の公比を求める.行列との関係は\sectionref*{節：漸化式と行列}で扱う. \booknavref{節：特性方程式型}{説明へ.}
  \item[\hypertarget{term:superposition}{重ね合わせの原理}] 同次の線形漸化式の解を定数倍して足したものも,また解になる性質. \booknavref{節：答えの形を見通す}{説明へ.}
  \item[\hypertarget{term:combination}{線形結合}] 数列やベクトルに定数を掛けて足したもの. \booknavref{節：答えの形を見通す}{説明へ.}
  \item[\hypertarget{term:basis}{基底}] どの対象もその組の線形結合としてただ1通りに表せるような組. \booknavref{節：二項係数で多項式を表す}{説明へ.}
\end{description}
\paragraphbf{行列とベクトル}
\begin{description}[style=nextline,leftmargin=0pt,labelsep=0pt,itemsep=5pt]
  \item[\hypertarget{term:matrix}{行列}] 数を長方形に並べたもの.連立漸化式の係数をまとめて表すのに使う. \booknavref{節：行列の基本}{説明へ.}
  \item[\hypertarget{term:column-vector}{列ベクトル}] 複数の数を縦に並べたもの. \booknavref{節：行列の基本}{説明へ.}
  \item[\hypertarget{term:row-vector}{行ベクトル}] 複数の数を横に並べたもの. \booknavref{節：行列の基本}{説明へ.}
  \item[\hypertarget{term:identity}{単位行列}] 掛けても元のベクトルや行列を変えない行列. \booknavref{節：行列の基本}{説明へ.}
  \item[\hypertarget{term:eigenvector}{固有ベクトル}] 行列を掛けると,元のベクトルの定数倍になる,0でないベクトル. \booknavref{節：行列の基本}{説明へ.}
  \item[\hypertarget{term:left-eigen}{左固有ベクトル}] 行列を右から掛けると,自身の定数倍になる,0でない行ベクトル. \booknavref{節：行列の基本}{説明へ.}
  \item[\hypertarget{term:eigenvalue}{固有値}] 固有ベクトルに行列を掛けたとき,何倍になるかを表す数.固有値の方程式$\det(\lambda I-A)=0$の解として求める. \booknavref{節：行列の基本}{説明へ.}
  \item[\hypertarget{term:transpose}{転置}] 行と列を入れ替える操作. \booknavref{節：行列の基本}{説明へ.}
\end{description}
\paragraphbf{角度と複素数}
\begin{description}[style=nextline,leftmargin=0pt,labelsep=0pt,itemsep=5pt]
  \item[\hypertarget{term:trigonometric}{三角関数}] 角度を,円上の位置や辺の比と結び付ける関数. \booknavref{節：三角関数}{説明へ.}
  \item[\hypertarget{term:addition}{加法定理}] 角度の和を,元の角度の三角関数で表す公式. \booknavref{節：三角関数}{説明へ.}
  \item[\hypertarget{term:complex-plane}{複素平面}] 複素数を平面上の点として表した平面. \booknavref{節：複素数を平面に表す}{説明へ.}
  \item[\hypertarget{term:polar}{極形式}] 複素数を,絶対値と角度によって表す形. \booknavref{節：複素数を平面に表す}{説明へ.}
  \item[\hypertarget{term:conjugate}{共役}] 複素数の虚部の符号を反転させる関係. \booknavref{節：複素数を平面に表す}{説明へ.}
  \item[\hypertarget{term:argument}{偏角}] 複素平面で,正の実軸から測る角度.角度の範囲の取り方を伴う. \booknavref{節：複素数を平面に表す}{説明へ.}
  \item[\hypertarget{term:de-moivre}{ド・モアブルの定理}] 複素数の整数乗で,絶対値と角度がどう変わるかを表す関係. \booknavref{節：ド・モアブルの定理}{説明へ.}
  \item[\hypertarget{term:euler}{オイラーの公式}] 複素指数関数と正弦・余弦を結ぶ公式. \booknavref{節：オイラーの公式}{説明へ.}
\end{description}
\paragraphbf{数列の広がり}
\begin{description}[style=nextline,leftmargin=0pt,labelsep=0pt,itemsep=5pt]
  \item[\hypertarget{term:invariant}{不変量}] 添字を1つ進めても値が変わらない量. \booknavref{節：変わらない量を使う}{説明へ.}
  \item[\hypertarget{term:newton}{ニュートン法}] 曲線の接線と$x$軸の交点を,方程式の解の次の近似値にする方法. \booknavref{節：ニュートン法}{説明へ.}
  \item[\hypertarget{term:limit}{極限}] 番号や変数を近づけるときに,値が近づく先を調べる考え方. \booknavref{節：極限}{説明へ.}
  \item[\hypertarget{term:one-sided}{片側極限}] 変数を片側から近づけたときの極限. \booknavref{節：片側極限}{説明へ.}
  \item[\hypertarget{term:left-limit}{左側極限}] 小さい側から変数を近づけたときの極限. \booknavref{節：片側極限}{説明へ.}
  \item[\hypertarget{term:right-limit}{右側極限}] 大きい側から変数を近づけたときの極限. \booknavref{節：片側極限}{説明へ.}
  \item[\hypertarget{term:pell}{ペル方程式}] 平方数でない正の整数$D$に対し,\,$x^2-Dy^2=1$の整数解を求める方程式. \booknavref{節：整数から作る近似分数}{説明へ.}
  \item[\hypertarget{term:continued}{連分数}] 分母の中に分数が続く表示.近似分数を作る漸化式と関係が深い. \booknavref{節：整数から作る近似分数}{説明へ.}
  \item[\hypertarget{term:difference-operator}{差分}] 添字を1つずらした値から元の値を引く操作.階差を作る操作を記号$\Delta$で表したもの. \booknavref{節：和と差から微積分へ}{説明へ.}
  \item[\hypertarget{term:generating}{母関数}] 数列の項を係数に並べ,数列全体を1つの式として扱う道具. \booknavref{節：数列をまとめて見る}{説明へ.}
  \item[\hypertarget{term:convolution}{畳み込み}] 添字の和を合わせた項どうしの積を足す操作.母関数の積に現れる. \booknavref{節：数列をまとめて見る}{説明へ.}
  \item[\hypertarget{term:chebyshev}{チェビシェフ多項式}] 余弦の整数倍の公式を,多項式の列として表したもの. \booknavref{節：角度の整数倍}{説明へ.}
  \item[\hypertarget{term:legendre}{ルジャンドル多項式}] 三項漸化式で作られる多項式列.本文では直交性も紹介する. \booknavref{節：漸化式と行列}{説明へ.}
  \item[\hypertarget{term:interpolation}{ラグランジュ補間}] 有限個の点を通る多項式を作る方法. \booknavref{節：漸化式と行列}{説明へ.}
\end{description}
